\documentclass{article}
\usepackage{amsmath,amssymb,amsthm}
\begin{document}

\textbf{Subproblem 1 --- Lower bound and rigidity in the equality case.}

\medskip
\textbf{Context.}
Let $S\subset\mathbb{R}^2$ be a finite set of $n$ points in \emph{convex position} (every
point of $S$ is a vertex of $\operatorname{conv}(S)$).
Let $T\subset\mathbb{R}^2$ be a finite set of translation vectors, and let
$\varepsilon:T\to\{+1,-1\}$.
Define
\[
  F(p) = \sum_{\substack{t\in T\\ p-t\in S}} \varepsilon(t), \qquad p\in\mathbb{R}^2.
\]
Write $P=\operatorname{conv}(S)$ and $Q=\operatorname{conv}(T)$.
A point $x\in S+T$ is called \emph{lonely} if there is a unique $t\in T$ with
$x-t\in S$.

\medskip
\textbf{Statement.}  Prove each of the following.
\begin{enumerate}
  \item \emph{(Lonely points exist.)} For every $s\in S$ there exists $t_s\in T$ such
        that $\ell_s := s + t_s$ is lonely.
  \item \emph{(Distinctness.)} The $n$ points $\ell_{s_1},\ldots,\ell_{s_n}$ are
        pairwise distinct.
  \item \emph{(Non-vanishing.)} $F(\ell_s)=\varepsilon(t_s)\neq 0$ for every $s\in S$.
  \item \emph{(Lower bound.)} $|\{p\in\mathbb{R}^2:F(p)\neq 0\}|\ge n$.
  \item \emph{(Rigidity.)} If $|\{p:F(p)\neq 0\}|=n$, then
        $\{p:F(p)\neq 0\}=\{\ell_s:s\in S\}$ and every point in $S+T$
        not of the form $\ell_s$ satisfies $F(p)=0$.
  \item \emph{(Vertices survive.)} If $|\{p:F(p)\neq 0\}|=n$, then every vertex of
        $P+Q$ belongs to $\{p:F(p)\neq 0\}$.
\end{enumerate}

\medskip
\textbf{Hint.}
For each $s\in S$, choose a linear functional $\phi_s:\mathbb{R}^2\to\mathbb{R}$
that is uniquely maximized at $s$ over $S$ (possible since $s$ is a vertex of
$\operatorname{conv}(S)$).  Maximize $\phi_s$ over $S+T$: the maximizer is
$s+t_s$ where $t_s=\arg\max_{t\in T}\phi_s(t)$.  If $t_s$ is not the unique
maximizer of $\phi_s$ over $T$, note that the argument still gives a point
uniquely covered by $S$: if $\ell_s = s+t_s \in S+t'$ for some $t'\neq t_s$,
then $\phi_s(\ell_s-t')=\phi_s(s+(t_s-t'))\le \phi_s(s)=\phi_s(\ell_s)-\phi_s(t_s)$
unless $\ell_s-t'=s$, i.e.\ $t'=t_s$.  For~(6), expose each vertex of $P+Q$
by a supporting linear functional and observe that the maximizer in $S+T$
is uniquely represented.

\end{document}
